\section{Cut reactions with graph diffusion}
\label{sec:spatial-diffusion}

Let $G=(V,E,c)$ be a connected weighted graph on $n\geq2$ vertices,
with symmetric nonnegative edge weights $c_{xy}$.  Its random-walk
Laplacian is
\[
 (\Delta_G a)(x)=\sum_yc_{xy}[a(x)-a(y)],
\]
and $\lambda_G>0$ denotes its smallest eigenvalue on the mean-zero
subspace.  On the nonempty subsets of $V$, define
\begin{equation}\label{eq:spatial-generator}
 \begin{split}
 (\mathcal L_{\beta,D,G}f)(S)
 ={}&\beta\sum_{x\notin S}[f(S\cup\{x\})-f(S)]\\
 &+\frac{|S|-1}{n}\sum_{x\in S}[f(S\setminus\{x\})-f(S)]\\
 &+D\sum_{\{x,y\}\in E}c_{xy}[f(S^{xy})-f(S)],
 \qquad \beta>0,\quad D\geq0,
 \end{split}
\end{equation}
where $S^{xy}$ exchanges the occupancy of $x$ and $y$.  The deletion
sum is zero at $|S|=1$.  The first two terms are the cut-reaction
generator, and the last term is conservative exclusion along $G$.
The reversible probability measure remains
\begin{equation}\label{eq:spatial-reversible-law}
 \pi_\beta(S)=Z^{-1}\frac{(\beta n)^{|S|}}{(|S|-1)!}.
\end{equation}
Indeed, its adjacent-cardinality weights satisfy detailed balance for
the reaction rates, and it is invariant under every coordinate
permutation.

The cardinality is the birth--death chain with rates
$\beta(n-k)$ and $k(k-1)/n$; write $\gamma_Q(n,\beta)$ for its gap.
The degree-one radial reaction operator acts on $k=1,\ldots,n-1$ as
\begin{equation}\label{eq:spatial-degree-one}
 \begin{split}
 (A_1h)(k)
 ={}&\beta(n-k-1)[h(k+1)-h(k)]\\
 &+\frac{(k-1)^2}{n}[h(k-1)-h(k)]
   -\left(\beta+\frac{k-1}{n}\right)h(k).
 \end{split}
\end{equation}
Let $\eta_1(n,\beta)$ be the smallest eigenvalue of $-A_1$.

\begin{theorem}[Exact gap with graph diffusion]
\label{thm:spatial-exact-gap}\label{thm:spatial-gap}
For every $n,\beta,D,G$ as above, the spectral gap of
\eqref{eq:spatial-generator} is
\begin{equation}\label{eq:spatial-exact-gap}
 \gamma_{\mathrm{sp}}(n,\beta,D,G)
 =\min\{\gamma_Q(n,\beta),\ \eta_1(n,\beta)+D\lambda_G\}.
\end{equation}
\end{theorem}

\begin{proof}
Use the Boolean harmonic decomposition of the cut chain.  In degree
$\ell$, functions are sums of products $h(|S|)H(S)$, where $H$ belongs
to the degree-$\ell$ harmonic space and $h$ is supported on the
allowable cardinalities.  The reaction generator acts on this sector
as $A_\ell\otimes I$, where $A_\ell=\mathcal L_\ell$.  Every coordinate permutation preserves harmonic
degree and commutes with the canonical lifts between slices.  Hence
the graph-swap generator acts as $I\otimes\mathcal S_{G,\ell}$, with
the same harmonic action at every cardinality.  The slice inner
products on the harmonic factor differ only by a scalar, so these
are tensor decompositions of the reversible Hilbert spaces as well.
Consequently the generator on the degree-$\ell$ sector is
\begin{equation}\label{eq:spatial-tensor-sum}
 A_\ell\otimes I+I\otimes D\mathcal S_{G,\ell}.
\end{equation}
The negative of the graph-swap operator is nonnegative, since it is
the restriction of a reversible Markov generator with its sign
changed.

In degree zero, swaps act trivially and the first positive eigenvalue
is $\gamma_Q$.  In every degree $\ell\geq2$,
Proposition~\ref{prop:higher-exact} gives
$\inf\operatorname{Spec}(-A_\ell)\geq\gamma_Q$; adding the graph form
cannot lower this bound.

For degree one write
\[
 H_a(S)=\sum_{x\in S}a(x),\qquad \sum_{x\in V}a(x)=0.
\]
A direct calculation gives
$\mathcal S_G H_a=-H_{\Delta_Ga}$.  Thus the degree-one eigenvalues
are sums of eigenvalues of $-A_1$ and of $D\Delta_G$ on mean-zero
vectors.  Their smallest value is $\eta_1+D\lambda_G$.  Taking the
minimum over all degrees proves~\eqref{eq:spatial-exact-gap}.
\end{proof}

Only the one-particle graph action is required in this proof.  Higher
harmonic sectors already lie above the count gap before diffusion is
added.

\begin{proposition}[An additive-error formula for the degree-one gap]
\label{prop:spatial-degree-one-sharp}
Set
\[
 d=\sqrt{\beta^2+4\beta},\qquad
 c_1=\frac{\beta+d}{2},\qquad a=nc_1.
\]
Then
\begin{equation}\label{eq:spatial-degree-one-endpoint-bracket}
 c_1-\frac{2a}{n(a+2)}
 \leq\eta_1(n,\beta)
 \leq c_1-\frac{2a^2}{n(n+a-2)(n+a)}.
\end{equation}
In particular,
\begin{equation}\label{eq:spatial-degree-one-additive-error}
 c_1-\frac2n\leq\eta_1(n,\beta)\leq c_1.
\end{equation}
\end{proposition}

\begin{proof}
The number $a$ is positive and satisfies $a^2=\beta n(n+a)$.
For $h(k)=1/(k+a)$ the quotient in
\eqref{eq:spatial-reciprocal-quotient} simplifies to
\begin{equation}\label{eq:spatial-sharp-reciprocal-quotient}
 -\frac{A_1h(k)}{h(k)}
 =\frac an-\frac{2a^2}{n(k+a-1)(k+a+1)}.
\end{equation}
Every denominator is positive, even when $a<1$, since $k\geq1$.
The quotient increases with $k$.  Its values at $k=1$ and $k=n-1$
give~\eqref{eq:spatial-degree-one-endpoint-bracket} by positive-vector
comparison.  The simpler bracket follows from $a/(a+2)\leq1$.
\end{proof}

\begin{corollary}[Leading constants with graph diffusion]
\label{cor:spatial-leading-constants}
With $d$ and $c_1$ as above, put
\[
 m=\min\{d,c_1+D\lambda_G\}.
\]
The count-chain bracket $d-5/n\leq\gamma_Q\leq d+4/n$ from
Theorem~\ref{thm:sharp-gap} and
Theorem~\ref{thm:spatial-exact-gap} imply
\begin{equation}\label{eq:spatial-leading-constants}
 m-\frac5n\leq\gamma_{\mathrm{sp}}(n,\beta,D,G)
 \leq m+\frac4n.
\end{equation}
In particular, if $n^2\beta\longrightarrow\infty$, the ratio of
$\gamma_{\mathrm{sp}}$ to $m$ tends to one, uniformly in the graph
and in $D\geq0$.
\end{corollary}

\begin{proof}
Combine Theorem~\ref{thm:sharp-gap} and the degree-one additive-error
bracket inside the minimum in
\eqref{eq:spatial-exact-gap}.  For the final assertion, observe that
$d\geq2\sqrt\beta$ and $c_1\geq\sqrt\beta$, so
$nm\geq n\sqrt\beta\longrightarrow\infty$.
\end{proof}

\begin{proposition}[The degree-one scale]
\label{prop:spatial-degree-one-scale}
For every $n\geq2$ and $\beta>0$, set
\[
 u=\beta+\min\{n\beta,\sqrt\beta\}.
\]
Then
\begin{equation}\label{eq:spatial-degree-one-scale}
 \frac u8\leq\eta_1(n,\beta)\leq u.
\end{equation}
More precisely, with $c_0=1-2^{-1/2}$,
\begin{equation}\label{eq:spatial-degree-one-auxiliary}
 \max\left\{\beta,\min\left(c_0n\beta,\frac1{2n}\right)\right\}
 \leq\eta_1
 \leq\min\left\{\frac{n\beta}{2},\beta+\sqrt\beta\right\}.
\end{equation}
Also,
\begin{equation}\label{eq:spatial-degree-one-below-count}
 \eta_1\leq\gamma_Q-\frac2n.
\end{equation}
In particular, the gap of the cut-reaction chain without graph
diffusion is exactly $\eta_1$.
\end{proposition}

\begin{proof}
The positive test function $h(k)=1/k$ gives the exact quotient
\[
 -\frac{A_1h(k)}{h(k)}=\frac{\beta n}{k+1}\leq\frac{n\beta}{2},
\]
which proves the first upper bound.  For the second, put $j=k-1$ and
consider the conservative part of $A_1$.  Its birth rate is
$\beta(n-2-j)$, its death rate is $j^2/n$, and its killing is
$\beta+j/n$.  If $\nu$ is the invariant law of this finite
conservative chain, stationarity applied to $j$ yields
\[
 \mathbb E_\nu j^2
 =\beta n(n-2-\mathbb E_\nu j)\leq\beta n^2.
\]
The constant-function Rayleigh quotient, followed by Cauchy--Schwarz,
therefore gives $\eta_1\leq\beta+\sqrt\beta$.  When $n=2$ the
conservative chain has the single state $j=0$, so the same argument
applies.

The killing term gives $\eta_1\geq\beta$.  For another lower bound
use $h(k)=k^{-1/2}$.  At $k=1$ its quotient is
\[
 \beta+\beta(n-2)(1-2^{-1/2})\geq c_0n\beta.
\]
For $k\geq2$, discard the nonnegative birth term and the term $\beta$
to obtain
\[
 -\frac{A_1h(k)}{h(k)}
 \geq\frac{k-1}{n}
 \left[1-(k-1)\left(\sqrt{\frac{k}{k-1}}-1\right)\right]
 \geq\frac1{2n}.
\]
Here
$(k-1)(\sqrt{k/(k-1)}-1)=1/(\sqrt{k/(k-1)}+1)\leq1/2$.
Positive-vector comparison proves
\eqref{eq:spatial-degree-one-auxiliary}.

It remains to obtain a square-root lower bound in the intermediate
range.  If $n\geq8$ and $4/n^2\leq\beta\leq1/16$, put
$a=n\sqrt\beta$, so that $2\leq a\leq n/4$, and use
$h(k)=1/(k+a)$.  Direct cancellation gives
\begin{equation}\label{eq:spatial-reciprocal-quotient}
 -\frac{A_1h(k)}{h(k)}
 =\beta+\frac{\beta(n-k-1)}{k+a+1}
       +\frac{a(k-1)}{n(k+a-1)}.
\end{equation}
If $k\leq a$, then $n-k-1\geq n/2$ and $k+a+1\leq5a/2$,
so the second term is at least $\sqrt\beta/5$.  If $k\geq a$,
then $k-1\geq k/2$ and $k+a-1\leq2k$, so the final term is at
least $\sqrt\beta/4$.  Thus $\eta_1\geq\sqrt\beta/5$ in this
range.

The estimates patch with a universal constant.  For
$\beta\leq n^{-2}$, the lower bound $c_0n\beta$ is at least
$(n+1)\beta/8=u/8$.  For $n\leq7$ and $\beta\geq n^{-2}$,
$u\leq8\beta$.  If $n\geq8$ and
$n^{-2}\leq\beta\leq4n^{-2}$, the auxiliary bound is at least
$c_0/n\geq\sqrt\beta/7$; combine this with $\eta_1\geq\beta$.
In the intermediate range just treated the two lower bounds give
$u\leq6\eta_1$.  Finally, for $\beta\geq1/16$, $u\leq5\beta$.
These cases prove the lower bound in~\eqref{eq:spatial-degree-one-scale}.
The two upper bounds in~\eqref{eq:spatial-degree-one-auxiliary}
immediately give $\eta_1\leq u$.

Finally, let $h$ be the positive ground state of the count-derivative
operator $A_{n,\beta}$.  Its forward ratios are at least $1/2$ by
Lemma~\ref{lem:critical-ground-ratios} (or by the weaker ratio bound
used in Proposition~\ref{prop:higher-exact}).  Direct subtraction gives
\[
 \frac{(A_{n,\beta}-A_1)h(k)}{h(k)}
 =\frac{(k-1)h(k-1)/h(k)-2k}{n}\leq-\frac2n.
\]
At $k=1$ the backward term is zero.  Testing $-A_1$ with $h$ proves
\eqref{eq:spatial-degree-one-below-count}, and
Theorem~\ref{thm:spatial-exact-gap} gives the last assertion.
\end{proof}

\begin{corollary}[A uniform spatial phase diagram]
\label{cor:spatial-uniform-gap}
Put $s=\beta+\sqrt\beta+n^{-1}$ and
$u=\beta+\min\{n\beta,\sqrt\beta\}$.  Then
\begin{equation}\label{eq:spatial-uniform-gap}
 \frac18\min\{s,u+D\lambda_G\}
 \leq\gamma_{\mathrm{sp}}(n,\beta,D,G)
 \leq4\min\{s,u+D\lambda_G\}.
\end{equation}
In particular, for $0<\beta\leq n^{-2}$,
\begin{equation}\label{eq:spatial-rare-gap}
 \gamma_{\mathrm{sp}}(n,\beta,D,G)
 \asymp\min\{n^{-1},n\beta+D\lambda_G\},
\end{equation}
with absolute comparison constants.
\end{corollary}

\begin{proof}
Combine Theorem~\ref{thm:spatial-exact-gap} and
Proposition~\ref{prop:spatial-degree-one-scale} with the count-chain
bounds $s/6\leq\gamma_Q\leq4s$.  In the stated rare-fragmentation
range, $s\asymp n^{-1}$ and $u\asymp n\beta$.
\end{proof}

For a unit-rate cycle on $n\geq3$ vertices,
$\lambda_G=4\sin^2(\pi/n)$.  For the discrete torus of side
$L\geq3$ in dimension $d$, with unit rates along its edges,
$n=L^d$ and $\lambda_G=4\sin^2(\pi/L)$.  Thus
\eqref{eq:spatial-rare-gap} becomes
\[
 \gamma_{\mathrm{sp}}
 \asymp\min\{n^{-1},n\beta+D/L^2\}
\]
on a torus.  This identifies the competition between count relaxation,
reactions that renew a dilute cut configuration, and graph diffusion.

On a cycle, adding or deleting a cut splits or merges neighboring
intervals, and exchanging adjacent occupied and unoccupied vertices
moves an interval boundary.  On a general graph,
\eqref{eq:spatial-generator} is an exclusion process with
cardinality-dependent creation and deletion.  Its reaction rates
retain the global count dependence and its equilibrium law remains
exchangeable.  These features permit the tensor decomposition and
specify the scope of the spatial result.

\subsection{Inhomogeneous reversible reactions}

The exact formula also controls reaction rates that depend on the
positions and surrounding configuration, through a comparison of
reversible conductances.  For each undirected reaction edge
$e=\{S,S\cup\{x\}\}$ with $S\neq\varnothing$ and $x\notin S$,
choose a factor $a_e$ satisfying
\[
 0<a_-\leq a_e\leq a_+<\infty.
\]
Multiply both the addition rate from $S$ to $S\cup\{x\}$ and the
reverse deletion rate by this same factor.  Keep the graph-swap
generator and its diffusion coefficient $D$ unchanged.  Write
$\gamma_{\mathrm{inhom}}$ for the resulting gap.

\begin{corollary}[Bounded inhomogeneous reaction rates]
\label{cor:spatial-inhomogeneous}
The modified chain is reversible for~\eqref{eq:spatial-reversible-law}
and satisfies
\begin{equation}\label{eq:spatial-inhomogeneous-gap}
 \begin{split}
 \min\{a_-\gamma_Q,\ a_-\eta_1+D\lambda_G\}
 &\leq\gamma_{\mathrm{inhom}}\\
 &\leq\min\{a_+\gamma_Q,\ a_+\eta_1+D\lambda_G\}.
 \end{split}
\end{equation}
The factors may depend arbitrarily on the edge and its configuration;
the modified reaction generator need not commute with permutations.
\end{corollary}

\begin{proof}
Multiplication by the same factor in both directions preserves
detailed balance at each reaction edge.  Let $\mathcal E_B$ be the
reaction form and $\mathcal E_G$ the unit-diffusion graph-swap form,
both with respect to $\pi_\beta$.  The modified form satisfies
\[
 a_-\mathcal E_B(f,f)+D\mathcal E_G(f,f)
 \leq\mathcal E_{\mathrm{inhom}}(f,f)
 \leq a_+\mathcal E_B(f,f)+D\mathcal E_G(f,f).
\]
For any constant $a>0$, Theorem~\ref{thm:spatial-exact-gap}, with
diffusion coefficient $D/a$ and then a change of time by $a$, gives
the gap of $aB+D\mathcal S_G$ as
$\min\{a\gamma_Q,a\eta_1+D\lambda_G\}$.  All three forms have the
same variance denominator.  Taking their infima proves
\eqref{eq:spatial-inhomogeneous-gap}.
\end{proof}

For example, place the vertices in a cyclic order and let adding a
cut at $x$ divide an interval into two positive lengths $i$ and $j$.
For $|\varepsilon|<1/2$, set
\begin{equation}\label{eq:spatial-circle-inhomogeneity}
 a_{\{S,S\cup\{x\}\}}
 =1+\varepsilon\cos\left(\frac{2\pi x}{n}\right)
   +\varepsilon\frac{(i-j)^2}{(i+j)^2}.
\end{equation}
This combines a spatially varying environment with a dependence on
the neighboring interval lengths.  The factor is unchanged when
$i$ and $j$ are interchanged, and the same split determines it for
the reverse merger.  It satisfies
$1-2|\varepsilon|\leq a_e\leq1+2|\varepsilon|$, so
Corollary~\ref{cor:spatial-inhomogeneous} applies with these constants.
The corollary covers bounded reversible modifications of the stated
reaction edges; its proof does not extend the exact formula to
arbitrary coagulation or fragmentation kernels.
